\documentclass[a4paper,11pt]{ltjarticle}
\usepackage{amsmath, amssymb, bm}
\usepackage{physics}
\usepackage[margin=2cm]{geometry}

\title{時間依存ゴースト・グッツウィラー近似 (TD-gGA) の定式化と導出}
\author{}
\date{}

\begin{document}
\maketitle

\section{平衡状態のgGAラグランジアン (Lanataの一般論)}
平衡状態におけるゴースト・グッツウィラー近似(gGA)では、系の基底状態エネルギーを最小化するために、補助空間の準粒子状態 $\ket{\Psi_0}$ と各サイト $i$ の埋め込み状態 $\ket{\Phi_i}$ を最適化します。この変分問題は、以下のラグランジアン $\mathcal{L}_{\text{eq}}$ の停留点（サドルポイント）を求める問題に帰着します。

\begin{align}
\mathcal{L}_{\text{eq}} &= \mel{\Psi_0}{\hat{H}_{qp}[\mathcal{R}, \Lambda]}{\Psi_0} + E(1 - \braket{\Psi_0}{\Psi_0}) \notag \\
&\quad + \sum_{i=1}^{\mathcal{N}} \left[ \mel{\Phi_i}{\hat{H}_{emb}^i[\mathcal{D}_i, \Lambda_i^c]}{\Phi_i} + E_i^c(1 - \braket{\Phi_i}{\Phi_i}) \right] \notag \\
&\quad - \sum_{i=1}^{\mathcal{N}} \left[ \sum_{a,b} (\Lambda_{i,ab} + \Lambda_{i,ab}^c)\Delta_{i,ab} \right. \notag \\
&\quad \left. + \sum_{c,a,\alpha} \left( \mathcal{D}_{i,a\alpha} \mathcal{R}_{i,c\alpha} [\Delta_i(1-\Delta_i)]_{ca}^{\frac{1}{2}} + c.c. \right) \right]
\label{eq:L_eq}
\end{align}
ここで、$E, E_i^c$ は状態のノルム（規格化）を保存するためのラグランジュ乗数であり、$\Lambda, \Lambda^c, \mathcal{D}$ はGutzwiller制約条件を満たすためのラグランジュ乗数です。

\section{ディラック・フレンケルの変分原理と時間依存ラグランジアン}
系が非平衡ダイナミクスを示す場合（例：相互作用クエンチや外部電磁場がある場合）、時間依存シュレーディンガー方程式の厳密な時間発展を変分空間内に制限して近似的に解く必要があります。これにはディラック・フレンケル（Dirac-Frenkel）の変分原理を用います。

作用積分 $\mathcal{S} = \int \mathcal{L}(t) dt$ の被積分関数となる時間依存ラグランジアン $\mathcal{L}(t)$ を構築するため、式(\ref{eq:L_eq})に対して以下の置き換えを行います。
\begin{enumerate}
    \item エネルギー期待値 $\mel{\Psi}{\hat{H}}{\Psi}$ を、時間発展演算子を含む $\mel{\Psi}{i\partial_t - \hat{H}}{\Psi}$ に置き換える。
    \item ユニタリな時間発展 $i\partial_t$ は状態のノルムを自動的に保存するため、規格化のためのラグランジュ乗数項 ($E, E_i^c$) を削除する。
    \item 各変分パラメータと状態ベクトルを時間 $t$ の関数とする。
\end{enumerate}

これにより、TD-gGAの時間依存ラグランジアンは次のように定義されます。
\begin{align}
\mathcal{L}(t) &= \mel{\Psi_0(t)}{i\partial_t - \hat{H}_{qp}(t)}{\Psi_0(t)} + \sum_{i=1}^{\mathcal{N}} \mel{\Phi_i(t)}{i\partial_t - \hat{H}_{emb}^i(t)}{\Phi_i(t)} \notag \\
&\quad + \sum_{i=1}^{\mathcal{N}} \left[ \sum_{a,b} (\Lambda_{i,ab}(t) + \Lambda_{i,ab}^c(t))\Delta_{i,ab}(t) \right. \notag \\
&\quad \left. + \sum_{c,a,\alpha} \left( \mathcal{D}_{i,a\alpha}(t) \mathcal{R}_{i,c\alpha}(t) [\Delta_i(t)(1-\Delta_i(t))]_{ca}^{\frac{1}{2}} + c.c. \right) \right]
\label{eq:L_td}
\end{align}

\section{運動方程式の導出（サドルポイント条件）}
作用積分 $\mathcal{S} = \int \mathcal{L}(t) dt$ を極値化するため、各変数に対する変分 $\delta \mathcal{S} = 0$ を計算します。

\subsection{状態ベクトルに関する変分}
$\bra{\Psi_0(t)}$ および $\bra{\Phi_i(t)}$ で変分をとると、それぞれ対応するハミルトニアンによる時間依存シュレーディンガー方程式が得られます。
\begin{align}
\frac{\delta \mathcal{S}}{\delta \bra{\Psi_0(t)}} = 0 \quad &\Rightarrow \quad i\partial_t \ket{\Psi_0(t)} = \hat{H}_{qp}(t) \ket{\Psi_0(t)} \label{eq:eom_psi} \\
\frac{\delta \mathcal{S}}{\delta \bra{\Phi_i(t)}} = 0 \quad &\Rightarrow \quad i\partial_t \ket{\Phi_i(t)} = \hat{H}_{emb}^i(t) \ket{\Phi_i(t)} \label{eq:eom_phi}
\end{align}

\subsection{ラグランジュ乗数に関する変分}
次に、制約条件のラグランジュ乗数 $\Lambda_{i,ab}(t), \Lambda_{i,ab}^c(t), \mathcal{D}_{i,a\alpha}(t)$ などで変分をとります。ハミルトニアン $\hat{H}_{qp}$ と $\hat{H}_{emb}^i$ はこれらの変数に対して線形であるため、得られる条件式は時間微分のない代数方程式となります。

\begin{align}
\frac{\delta \mathcal{S}}{\delta \Lambda_{i,ab}(t)} = 0 \quad &\Rightarrow \quad \Delta_{i,ab}(t) = \mel{\Psi_0(t)}{f_{ia}^\dagger f_{ib}}{\Psi_0(t)} \\
\frac{\delta \mathcal{S}}{\delta \Lambda_{i,ab}^c(t)} = 0 \quad &\Rightarrow \quad \Delta_{i,ab}(t) = \mel{\Phi_i(t)}{b_{ib}^\dagger b_{ia}}{\Phi_i(t)} \\
\frac{\delta \mathcal{S}}{\delta \mathcal{D}_{i,a\alpha}(t)} = 0 \quad &\Rightarrow \quad \mel{\Phi_i(t)}{c_{i\alpha}^\dagger b_{ia}}{\Phi_i(t)} = \sum_{c} [\Delta_i(t)(1-\Delta_i(t))]_{ac}^{\frac{1}{2}} \mathcal{R}_{i,c\alpha}(t)
\end{align}
驚くべきことに、これらの代数的な制約方程式は、各瞬間 $t$ において平衡状態のgGAと全く同じ形をしています。

\section{一般化されたTD-gGAのラグランジュ方程式群 (Lanata記法準拠)}
無限次元極限に依存しない、一般的な多軌道・有限格子系（サイト $i, j$、軌道 $\alpha, \beta$、ゴースト軌道 $a, b$）に対するTD-gGAのサドルポイント条件を記述する。これはLanataの平衡状態の方程式(89)-(95)をそのまま時間領域へ拡張したものである。

\subsection{1. 時間発展方程式（微分方程式）}
状態ベクトルの変分から、埋め込み状態と準粒子状態の時間発展方程式が導出される。

\paragraph{埋め込み状態の時間発展}
各サイト $i$ の局所的なゴースト自由度を含む埋め込み状態は、以下の微分方程式に従う。
\begin{align}
    i\partial_t \ket{\Phi_i(t)} &= \hat{H}_{emb}^i(t) \ket{\Phi_i(t)} \label{eq:td_phi_pure}
\end{align}

\paragraph{準粒子状態の時間発展}
系全体の補助フェルミオンからなるスレイター行列式（準粒子状態）は、以下の微分方程式に従う。
\begin{align}
    i\partial_t \ket{\Psi_0(t)} &= \hat{H}_{qp}(t) \ket{\Psi_0(t)} \label{eq:td_psi_pure}
\end{align}

\subsection{2. 自己無撞着条件（代数方程式）}
各時刻 $t$ において、変分パラメータは以下の4つの代数方程式をサイトごとに満たさなければならない。

\paragraph{条件1: $\Delta_i(t)$ の補助粒子からの定義}
\begin{align}
    [\Delta_i]_{ab}(t) &= \mel{\Psi_0(t)}{f_{ia}^\dagger f_{ib}}{\Psi_0(t)} = \rho_{ii, ab}(t) \label{eq:td_delta_psi}
\end{align}
※ 論文にあった $\int d\omega D_0(\omega)$ の積分は、一般の格子系では「系全体の密度行列 $\bm{\rho}(t)$ のサイト $i$ の対角ブロックを抽出する」というこの操作に置き換わる。

\paragraph{条件2: $\Delta_i(t)$ の埋め込み状態からの定義}
\begin{align}
    [\Delta_i]_{ab}(t) &= \mel{\Phi_i(t)}{b_{ib}^\dagger b_{ia}}{\Phi_i(t)} \label{eq:td_delta_phi}
\end{align}

\paragraph{条件3: $\mathcal{R}_i(t)$ の決定方程式}
\begin{align}
    \mel{\Phi_i(t)}{c_{i\alpha}^\dagger b_{ia}}{\Phi_i(t)} &= \sum_{c=1}^{B\nu_i} [\Delta_i(t)(1-\Delta_i(t))]^{\frac{1}{2}}_{ac} [\mathcal{R}_i]_{c\alpha}(t) \label{eq:td_R_gen}
\end{align}

\paragraph{条件4: $\Lambda_i^c(t)$ の決定方程式}
\begin{align}
    [\Lambda_i^c]_{ab}(t) &= - [\Lambda_i]_{ab}(t) - \sum_{c, \alpha} \frac{\partial}{\partial [\Delta_i]_{ab}} \left( [\Delta_i(t)(1-\Delta_i(t))]^{\frac{1}{2}}_{cb} [\mathcal{D}_i]_{b\alpha}(t) [\mathcal{R}_i]_{c\alpha}(t) + c.c. \right) \label{eq:td_lambda_gen}
\end{align}

\section{準粒子密度行列の運動方程式の導出（詳細）}
巨大なスレイター行列式 $\ket{\Psi_0(t)}$ のシュレーディンガー方程式による時間発展を、計算コストの低い1粒子密度行列 $n_{ab}(t)$ の運動方程式へと変換する数学的なプロセスを詳細に示す。

準粒子ハミルトニアン $\hat{H}_{qp}(t)$ は相互作用を含まない1体演算子であり、次のように一般化して書くことができる。
\begin{align}
    \hat{H}_{qp}(t) = \sum_{c,d} h_{cd}(t) f_c^\dagger f_d
\end{align}
ここで、$h_{cd}(t)$ はハミルトニアンの行列要素である。

\subsection{ステップ1：密度行列の時間微分}
1粒子密度行列の定義 $n_{ab}(t) = \mel{\Psi_0(t)}{f_a^\dagger f_b}{\Psi_0(t)}$ の両辺を時間 $t$ で微分する。ここに、シュレーディンガー方程式 $i\partial_t \ket{\Psi_0} = \hat{H}_{qp}\ket{\Psi_0}$ とそのエルミート共役 $-i\partial_t \bra{\Psi_0} = \bra{\Psi_0}\hat{H}_{qp}$ を代入する。
\begin{align}
    i\partial_t n_{ab}(t) &= i (\partial_t \bra{\Psi_0}) f_a^\dagger f_b \ket{\Psi_0} + i \bra{\Psi_0} f_a^\dagger f_b (\partial_t \ket{\Psi_0}) \notag \\
    &= -\bra{\Psi_0} \hat{H}_{qp} f_a^\dagger f_b \ket{\Psi_0} + \bra{\Psi_0} f_a^\dagger f_b \hat{H}_{qp} \ket{\Psi_0} \notag \\
    &= \mel{\Psi_0}{\left[ f_a^\dagger f_b, \hat{H}_{qp} \right]}{\Psi_0}
\end{align}
これにより、ハイゼンベルクの運動方程式と同等の交換子の期待値に帰着する。

\subsection{ステップ2：1体ハミルトニアンとの交換関係}
次に、フェルミオンの反交換関係 $\{f_x, f_y^\dagger\} = \delta_{xy}$ と、公式 $[AB, CD] = A\{B,C\}D - C\{A,D\}B$ を用いて、演算子の交換子を計算する。
\begin{align}
    \left[ f_a^\dagger f_b, f_c^\dagger f_d \right] &= f_a^\dagger \{f_b, f_c^\dagger\} f_d - f_c^\dagger \{f_a^\dagger, f_d\} f_b \notag \\
    &= \delta_{bc} f_a^\dagger f_d - \delta_{ad} f_c^\dagger f_b
\end{align}
これを $\hat{H}_{qp}(t)$ の定義に代入すると、クロネッカーのデルタによって和が簡約される。
\begin{align}
    \left[ f_a^\dagger f_b, \hat{H}_{qp}(t) \right] &= \sum_{c,d} h_{cd}(t) \left( \delta_{bc} f_a^\dagger f_d - \delta_{ad} f_c^\dagger f_b \right) \notag \\
    &= \sum_{d} h_{bd}(t) f_a^\dagger f_d - \sum_{c} h_{ca}(t) f_c^\dagger f_b
\end{align}

\subsection{ステップ3：方程式の閉包}
ステップ2で得られた演算子を再び $\bra{\Psi_0(t)} \dots \ket{\Psi_0(t)}$ で挟んで期待値をとる。これにより、再び密度行列 $n_{xy}$ の定義を用いることができる。
\begin{align}
    i\partial_t n_{ab}(t) &= \sum_{d} h_{bd}(t) \mel{\Psi_0}{f_a^\dagger f_d}{\Psi_0} - \sum_{c} h_{ca}(t) \mel{\Psi_0}{f_c^\dagger f_b}{\Psi_0} \notag \\
    &= \sum_{c} \left( h_{bc}(t) n_{ac}(t) - h_{ca}(t) n_{cb}(t) \right) \label{eq:eom_density_derivation}
\end{align}
（※2行目ではダミーインデックス $d$ を $c$ に統一した。）

式(\ref{eq:eom_density_derivation})は、行列の記法を用いれば $i\partial_t \bm{n} = [\bm{n}, \bm{h}^T]$ と等価である。gGAの枠組みにより、相互作用をすべて埋め込みハミルトニアン $\hat{H}_{emb}$ に分離し、$\hat{H}_{qp}$ を1体演算子に制限したことで、多体系のシュレーディンガー方程式を「密度行列のみで閉じた厳密な微分方程式」へと還元できたことが示された。

\section{アルゴリズムの実装に向けた変形（密度行列の時間発展）}
数値計算上、巨大な多体系スレイター行列式 $\ket{\Psi_0(t)}$ の時間発展（式\ref{eq:eom_psi}）を直接計算することは困難です。しかし、$\hat{H}_{qp}(t)$ は1体演算子であるため、1粒子密度行列 $n_{ab}(t) = \mel{\Psi_0(t)}{f_a^\dagger f_b}{\Psi_0(t)}$ の時間発展に完全に書き換えることができます。

ハイゼンベルクの運動方程式 $i\partial_t \hat{n}_{ab} = [\hat{n}_{ab}, \hat{H}_{qp}(t)]$ を用いることで、式(\ref{eq:eom_psi})は次のような密度行列の非線形常微分方程式に帰着します。
\begin{align}
i\partial_t n_{ab}(t) &= \sum_c \left[ h_{qp, bc}(t) n_{ac}(t) - h_{qp, ca}(t) n_{cb}(t) \right]
\label{eq:eom_density}
\end{align}
ここで、$h_{qp}(t)$ は準粒子ハミルトニアンの行列要素であり、行列 $\mathcal{R}(t)$ や $\Lambda(t)$ に依存します。

結果として、TD-gGAの方程式系は、埋め込み状態の時間発展（式\ref{eq:eom_phi}）、密度行列の時間発展（式\ref{eq:eom_density}）、および各時刻での代数的な自己無撞着条件をまとめた、ルンゲ・クッタ法で数値積分可能な巨大な1階常微分方程式 $\partial_t \bm{Y}(t) = \bm{F}(\bm{Y}(t))$ として定式化されます。

\end{document}